\section{Some explicit localization formulas}\label{appendix: some explicit formulas in type A}
This appendix is dedicated to spelling out some of our formulas for $GL_n$. We also recall the determinant line bundle and show that it only appears after Schubert completion.

\subsection{Formulas in \texorpdfstring{$GL_n$}{GLn}}
For the convenience of the reader, we record the description of the $K$-theory of $\Gr_{GL_n}$ together with its localization formulas explicitly in usual coordinates.


\subsubsection{Localization formulas in \texorpdfstring{$GL_n$}{GLn}.}\label{equation: localization formulas in GLn}
Consider $G = GL_n$. The connected components of $\Gr_{GL_n}$ are labeled by $\pi_1(GL_n) = \Z$.
We want to describe the localization map from functions on the corresponding union of affine blowups.
\begin{align*}
  \loc:  \prod_{\Z} \O(\eN_{\mu_{\infty}}(\eZ, T \times T \sslash W)) &\to \prod_{\lambda \in X_{\bullet}} \O(T \times \Gm^{\rot})
\end{align*}
Let $\lambda \in X_{\bullet}$ and denote $d_1, \dots, d_n$ the heights of the columns of its Young diagram. The localization $\loc_{\lambda}$ to the corresponding fixed point $t^{\lambda} = (t^{d_1}, \dots, t^{d_n})$ is given as follows. First project to the component labeled by $\overline{\lambda} = d_1 + d_2 + \dots + d_n \in \Z$. On generators of this copy, the localization map reads as
\begin{align*}
   \gamma_{\lambda}: \O(\eN_{\mu_{\infty}}(\eZ, T \times T \sslash W)) & \xrightarrow{} \O(T \times \Gm^{\rot}) \\
    t_i &\mapsto t_i \\
    q &\mapsto q \\
    s_i &\mapsto t_i q^{d_i} \\    
    e_j(s_1, \dots, s_n) &\mapsto e_j(t_1 q^{d_1}, \dots, t_n q^{d_n}) \\
    b_{j, \ell} = \tfrac{e_j(t^{\ell}_1, \dots, t^{\ell}_n) - e_j(s^{\ell}_1, \dots, s^{\ell}_n)}{1-q^{\ell}} &\mapsto \tfrac{e_j(t^{\ell}_1, \dots, t^{\ell}_n) - e_j(t_1 \cdot q^{d_1}, \dots, t_n \cdot q^{d_n})}{1-q^{\ell}}.
\end{align*}
Also note that these formulas work uniformly with respect to the connected components of $\Gr_{GL_n}$.



\subsubsection{Determinant line bundle \texorpdfstring{$\eL_{\det}$}{L det}.}
Let us first recall the construction of the determinant line bundle $\eL_{\det}$ on $\Gr_{GL_n}$.

\begin{recollection}
To a perfect complex $\eE \in \Perf(\eX)$ one can associate a graded line bundle $\det(\eL) \in \Pic(\eX)^{\Z}$. If we represent $\eE$ by an actual bounded complex of vector bundles on $\eX$, the determinant line bundle is given by an alternating product of the top wedge powers of the terms of $\eE$, while the grading records the alternating sum of their ranks. This descends to a map $K(\eX) \to \Pic(\eX)^{\Z}$. See \cite[Tag 0FJI]{Sta}, \cite[Proposition 5.3, \S 5, \S 12]{BS16}. 

In particular, to any class in $K^{0}(\eX)$ one can associate the isomorphism class of the corresponding line bundle; similarly in the topological setting.
\end{recollection}


\begin{construction}
We obtain the \textit{determinant line bundle} $\eL_{\det}$ on $\Gr_{GL_n}$ as follows. On each $\Gr_{\leq \mu}$, set
\begin{equation*}
    \eL_{\leq \mu, \det} := \det([\eE_{\leq \mu}] - [\eG_{\leq \mu}]).
\end{equation*}
These are compatible in $\mu$, and hence give a well-defined equivariant line bundle on the whole $\Gr_{GL_n}$. For each $\sigma \in \pi_1(G)$, it is known to be an ample generator of the Picard group $\Pic(\Gr^{\sigma}) = \Z$.

Also note that $\eL_{\det}$ is naturally $\lop GL_n \rtimes \Gm^{\rot}$-equivariant. It is not equivariant for the whole loop group $\lo G \rtimes \Gm^{\rot}$; this is corrected by working with respect to its central extension. 

For fixed $\mu$, $[\eL_{\det}] = \det(b_{1,1})$ is expressible in terms of $\lambda$-operations on $b_{1,1}$. With rational coefficients, the restriction of $[\eL_{\det}]$ to any $\Gr_{\leq \mu}$ may thus be expressed in terms of $b_{j, \ell}$. However, this expression is not uniform in $\mu$ -- the number of terms grows. (We will see shortly that the class $[\eL_{\det}] \in K^{\top, 0}_{T \times \Gm^{\rot}}(\Gr)$ can be seen only after the Schubert completion of our affine blowup.)
\end{construction}


 
\begin{claim}\label{equation: localization of determinant line bundle in GLn} \label{claim: localization of determinant line bundle in GLn}
The class of the determinant line bundle $\eL_{\det}$ inside the GKM description is
\begin{equation*}
   \loc_{\lambda}( [\eL_{\det}] ) = \prod_{k=1}^n t^{d_k}_k \cdot q^{\binom{d_k}{2}}, \qquad \forall \lambda \in X_{\bullet}.
\end{equation*}
\end{claim}
\begin{proof}
Directly follows from the localization formulas for $b_{1,1}$ in \S \ref{section: bezrukavnikov classes for GL_n from tautological bundles}. 
\end{proof}

\begin{remark}
    One can easily make the sanity check that the above class $[\eL_{\det}]$ indeed satisfies the GKM conditions. Namely, given $\lambda$ and $\lambda' = s_{\beta, m}\cdot \lambda$, where $\beta$ corresponds to the reflection switching $i$ and $j$, we can compute
    \begin{align*}
       \loc_{\lambda}( [\eL_{\det}]) - \loc_{\lambda'}( [\eL_{\det}] ) 
       &= \left( \prod_{k=1}^n t^{d_k}_k \cdot q^{\binom{d_k}{2}} \right) \cdot \left( 1 - t_jt_i^{-1} q^{d_j -d_i +1} \right)
    \end{align*}
 which is divisible by $(1 - {\beta} \cdot q^{m}) = (1 - t_jt_i^{-1} q^{d_j -d_i +1})$.   
\end{remark}


\begin{remark}
For general reductive $G$, the Picard group of each connected component of $\Gr_G$ is given by $\Z$. We denote its ample generator by $\eL_{\ver}$ and call it the \emph{Verlinde line bundle} (we reserve the name \emph{determinant line bundle} for $GL_n$). Pulling back $\eL^{GL_n}_{\det}$ along representations $G \to GL_n$ gives nonnegative powers of its ample generator $\eL_{\ver}$; these can be made explicit. See \cite[Lemma 4.2, Remark 4.3]{YZ09} for a discussion.  
\end{remark}

To describe the ample generator $\eL_{\ver}$ for a given component of the affine Grassmannian $\Gr_G$ of a reductive group $G$, it is sufficient to understand the case of $G=G^{\sc}$. Here, we have the following.
\begin{lemma}\label{lemma: localization of the verlinde line bundle}
Assume $G$ is simply connected.
The localization $\loc ([\eL_{\ver}])$ is given as follows. For any $\lambda \in X_{\bullet}$, we have
\begin{equation*}
    \loc_{\lambda} ([\eL_{\ver}]) = t_{\lambda} \cdot q^{\tfrac{1}{2} \langle \lambda^{\vee}, \lambda \rangle} \in \O(T \times \Gm^{\rot}).
\end{equation*}
\end{lemma}
\begin{proof}
Well-known, for example a special case of \cite[Proposition 7.6]{Yun10}. Note that \cite{Yun10} describes the localization formulas for the first Chern class of powers and twists of $\eL_{\ver}$ on the corresponding affine flag variety. By taking in his notation $\kappa=1$ and $\chi=1$, we get the pullback of the Verlinde line bundle from $\Gr_G$. The equivariant first Chern class of the localization of $\eL_{\ver}$ determines its equivariant $K$-theory class by exponentiation.
\end{proof}    




\subsubsection{Necessity of Schubert completion.}
The following observation well illustrates the necessity of Schubert completion.
\begin{lemma}\label{lemma: necessity of Schubert completion}
The class of the determinant line bundle does not lies in the functions on the uncompleted deformation
\begin{equation*}
    [\eL_{\det}] \notin \O(\eN_{\mu_{\infty}}(\eZ, T \times \Gm^{\rot} \times T \sslash W)).
\end{equation*}
\end{lemma}
\begin{proof}
For the sake of contradiction, suppose otherwise. Then we could express $[\eL_{\det}]$ as a polynomial $f$ over $k$ in finitely many elements from
\begin{equation}\label{equation: elements}
 \{ t_i, q, e_j, b_{j, \ell} \mid i = 1, \dots, n; j = 1, \dots, n; \ell \in \N \}   
\end{equation}
For any such fixed $f$, the localization map $\gamma = (\gamma_{\lambda} \mid \lambda \in \Z^d)$ has the following property: The powers of $q$ in the family of polynomials
\begin{equation*}
 \gamma_{\lambda}(f) \in \k[t_1^{\pm 1}, \dots, t^{\pm 1}_n, q^{\pm 1}]   
\end{equation*}
grow linearly in $\lambda = (d_1, \dots, d_n)$. Indeed, this is true for the elements from \eqref{equation: elements} by \S \ref{equation: localization formulas in GLn} and is preserved under polynomial combinations.

However, the powers of $q$ in the family
\begin{equation*}
 \gamma_{\lambda}([\eL_{\det}]) \in \k[t_1^{\pm 1}, \dots, t^{\pm 1}_n, q^{\pm 1}]   
\end{equation*}
grow quadratically in $\lambda = (d_1, \dots, d_n)$ by \eqref{equation: localization of determinant line bundle in GLn}.

Therefore, there exists a sufficiently big $\lambda$, such that $\gamma_{\lambda}([\eL_{\det}]) \neq \gamma_{\lambda}(f)$, a contradiction.
\end{proof}

\begin{remark}
The statement of Lemma \ref{lemma: necessity of Schubert completion} follows for any reductive group $G$. Indeed all the classes in question are base-changed from the $GL_n$-case; the statement about linear vs. quadratic growth of $q$-powers for varying $\lambda$ still holds -- the localizations may be also expressed purely Lie theoretically.
\end{remark}

In other words, $[\eL_{\det}]$ is not present in the uncompleted ring $\O(\eN_{\mu_{\infty}}(T \times T \sslash W))$. At the same time $[\eL_{\det}]$ is as nice a line bundle on $\Gr_G$ as one could hope for. This makes it challenging to give a purely $K$-theoretical interpretation of the decompleted ring $\O(\eN_{\mu_{\infty}}(T \times T \sslash W))$.


\subsubsection{\texorpdfstring{$\lambda$}{Lambda}-rings and affine Grassmannian of \texorpdfstring{$GL_{\infty}$}{GL infinity}.}


\begin{remark}
When $G=\GL_n$ and the coefficient ring $\k$ is a field of characteristic zero, it is possible to do with less topological generators. As complete topological $\lambda$-ring over the equivariant base, $K^{\top, 0}_{T \times \Gm^{\rot}}(\Gr; \C)$ is generated by the single tautological element $b_{1,1}$.

For $GL_n$, the representation ring $R(GL_n; \k) = \O(T \sslash W)$ is a quotient of the free $\lambda$-ring in one variable. Since we are working over a field of characteristic zero, Newton's formulas imply that we may further express any $\lambda$ operation (an elementary symmetric polynomial) purely in terms of Adams operations (power sums).
\end{remark}


\begin{remark}
Assume $G = GL_n$. In our presentation of $K$-theory of $\Gr_{GL_n}$, we have used the tautological class $b_{1,1}$ and $\psi$-operations (or $\lambda$-operations) on it. However, for fixed $n$, there are relations between these classes. 

This is closely related to the following phenomenon. The representation ring $R(GL_n) = \O(T \sslash W)$ is a quotient of the free $\lambda$-ring in one variable. The free $\lambda$-ring in one variable is given by putting all $R(GL_n)$ together for varying $n$.    


To have a uniform presentation for the $K$-theory, one could consider the ind-scheme
\begin{equation*}
    \Gr_{GL_{\infty}} := \colim_n \Gr_{GL_n}
\end{equation*}
equivariant with respect to the extended loop group of the ind-group scheme $GL_{\infty} := \colim_n GL_n$.
\end{remark}